\chapter{Unit 6: Development of Surfaces \& 3D Solid Editing}

\section{A. What You'll Learn}
Development of surfaces is the geometric unrolling or unfolding of 3D solid surfaces into a 2D flat plane pattern. In this unit, you will master:
\begin{enumerate}
    \item Surface development principles and industrial sheet metal applications (ducts, pipes, hoppers, funnels, chimneys).
    \item Governing rules: **Stretch-out Line** (full perimeter) and **True Length ($TL$)** enforcement.
    \item **Parallel Line Method** for Prisms and Cylinders.
    \item **Radial Line Method** for Pyramids and Cones ($\theta = \frac{r}{L} \times 360^\circ$).
    \item **Triangulation Method** for transition pieces.
    \item Development procedures for regular and truncated solids.
    \item AutoCAD 3D Solid Editing \& Boolean operations (`UNION`, `SUBTRACT`, `INTERSECT`).
    \item AutoCAD 3D Navigation (`3DORBIT`, Visual Styles) and 2D view extraction (`VIEWBASE`).
\end{enumerate}

---

\section{B. Principles of Surface Development}

\begin{definition}[Development of Surfaces]
The process of unfolding or unrolling the lateral surfaces of a 3D solid object onto a single 2D plane surface to produce a flat pattern. When cut from sheet metal and folded along bend lines, it forms the 3D solid.
\end{definition}

\subsection{Industrial Applications}
Sheet metal fabrication of air conditioning ducts, boiler shells, hoppers, funnels, chimneys, automobile body panels, and packaging boxes.

\subsection{Key Terminology \& Governing Rules}
\begin{itemize}
    \item \textbf{Stretch-out Line:} The full length or perimeter of the base laid out in a straight line.
    \item \textbf{True Length ($TL$):} Every line on the development pattern **MUST represent the true length** of the corresponding edge on the solid. If an edge is inclined in orthographic views, its true length must be determined before drawing development.
    \item \textbf{Seam Line Selection:} The opening/closing seam line is usually chosen along the **shortest edge** to minimize welding or soldering length.
\end{itemize}

---

\section{C. Methods of Surface Development}

\begin{table}[htbp]
\centering
\small
\begin{tabularx}{\textwidth}{l X X}
\toprule
\textbf{Method Name} & \textbf{Applicability (Solids)} & \textbf{Governing Geometric Principle} \\
\midrule
\textbf{Parallel Line Method} & Prisms and Cylinders & Lateral edges/generators are parallel to each other. Stretch-out line length $=$ Base Perimeter. \\
\addlinespace
\textbf{Radial Line Method} & Pyramids and Cones & Lateral edges slant towards a single apex. Development is a sector of a circle with Apex as center and Slant Edge ($TL$) as radius. \\
\addlinespace
\textbf{Triangulation Method} & Transition pieces (e.g., square-to-round adapters) & Surface is divided into a series of triangles laid out adjacent to one another. \\
\bottomrule
\end{tabularx}
\caption{Classification of Surface Development Methods.}
\label{tab:dev-methods}
\end{table}

---

\section{D. Development Procedures for Prisms \& Pyramids}

\subsection{1. Truncated Prisms (Parallel Line Method)}

\begin{keyequation}[Prism Stretch-out Line Formula]
\begin{equation}
    \text{Stretch-out Line Length } (L) = \text{Perimeter of Base}
\end{equation}
$$\text{Square Prism: } L = 4 \times s, \quad \text{Hexagonal Prism: } L = 6 \times s, \quad \text{Cylinder: } L = \pi D$$
\end{keyequation}

\subsubsection{Step-by-Step Procedure}
1. Draw Top View and Front View of the prism with section cutting plane.
2. Draw horizontal stretch-out line aligned with base of Front View equal to base perimeter.
3. Divide stretch-out line into segments corresponding to base side lengths ($s$).
4. Erect vertical lines (generators) from division points.
5. Project intersection points of cutting plane from Front View horizontally across to corresponding vertical lines on development.
6. Join projected points with straight lines to form upper boundary of development pattern.

\subsection{2. Truncated Pyramids \& Cones (Radial Line Method)}

\begin{keyequation}[Cone Development Sector Angle ($\theta$)]
For a cone with base radius $r$ and true slant height $L$:
\begin{equation}
    \theta = \left( \frac{r}{L} \right) \times 360^\circ
\end{equation}
\end{keyequation}

\subsubsection{Step-by-Step Procedure}
1. Draw Top and Front Views. Determine **True Length ($TL$)** of slant edge. (If slant edge in FV is not parallel to $XY$, rotate it to find $TL$).
2. Draw an arc with Apex $P$ as center and radius equal to $TL$.
3. Step off chord lengths along arc equal to base polygon side length $s$.
4. Join chord division points to Apex $P$ with straight lines.
5. **For Truncated Pyramids:** Mark cut points on the $TL$ slant edge in Front View. Transfer these distances to radial lines using concentric arcs centered at Apex $P$.

\begin{example}[Cone Development Sector Angle]
\textbf{Problem:} Calculate the sector angle ($\theta$) for developing a cone of base diameter $50\text{ mm}$ ($r = 25\text{ mm}$) and altitude height $60\text{ mm}$.\\[4pt]
\textbf{Solution Step 1: Calculate True Slant Height ($L$)}
$$L = \sqrt{r^2 + h^2} = \sqrt{25^2 + 60^2} = \sqrt{625 + 3600} = \sqrt{4225} = \mathbf{65\text{ mm}}$$

\noindent \textbf{Solution Step 2: Calculate Sector Angle ($\theta$)}
$$\theta = \left( \frac{r}{L} \right) \times 360^\circ = \left( \frac{25}{65} \right) \times 360^\circ = \mathbf{138.46^\circ}$$
\end{example}

---

\section{E. AutoCAD 3D Solid Editing \& Navigation}

\subsection{3D Boolean Operations}

\begin{autocadcmd}[UNION]
Merges two or more separate 3D solids or regions into a single composite solid.
\begin{itemize}
    \item \textbf{Command:} \texttt{UNION} or \texttt{UNI}
    \item \textbf{Workflow:} Type \texttt{UNI} $\to$ Select 3D solids $\to$ Press \texttt{Enter}.
\end{itemize}
\end{autocadcmd}

\begin{autocadcmd}[SUBTRACT]
Removes the volume of one 3D solid from another (essential for creating holes, slots, or truncated cuts).
\begin{itemize}
    \item \textbf{Command:} \texttt{SUBTRACT} or \texttt{SU}
    \item \textbf{Workflow:} Type \texttt{SU} $\to$ Select solid to KEEP $\to$ Press \texttt{Enter} $\to$ Select solid to REMOVE $\to$ Press \texttt{Enter}.
\end{itemize}
\end{autocadcmd}

\begin{autocadcmd}[INTERSECT]
Creates a 3D solid from the overlapping common volume of two or more solids.
\begin{itemize}
    \item \textbf{Command:} \texttt{INTERSECT} or \texttt{IN}
\end{itemize}
\end{autocadcmd}

\subsection{3D Navigation \& 2D View Extraction}
\begin{itemize}
    \item \textbf{\texttt{3DORBIT} (\texttt{3DO}):} Dynamically rotates 3D view in space (\texttt{Shift} + Middle Mouse Drag).
    \item \textbf{Visual Styles (\texttt{VS}):} Controls edge/shading display (\texttt{2D Wireframe}, \texttt{Conceptual}, \texttt{Realistic}, \texttt{Hidden}).
    \item \textbf{\texttt{VIEWBASE}:} Automatically generates linked 2D Orthographic Front, Top, Side, and Isometric views directly from 3D Model Space onto a Layout sheet.
\end{itemize}

---

\section{F. Chapter 6 Practice Problems \& MCQs}

\subsection{Practice Questions}

\begin{vivaqa}[Question 1 (Basic)]
Which development method is used for Prisms and Cylinders? What is the length of the stretch-out line?\\[4pt]
\textbf{Answer:} The **Parallel Line Method**. Length of stretch-out line equals the perimeter of the base ($L = \text{Perimeter}$).
\end{vivaqa}

\begin{vivaqa}[Question 2 (Moderate)]
Write the formula for the development sector angle ($\theta$) of a cone.\\[4pt]
\textbf{Answer:} $\theta = \left( \frac{r}{L} \right) \times 360^\circ$, where $r$ is base radius and $L$ is true slant height.
\end{vivaqa}

\begin{vivaqa}[Question 3 (Exam Level)]
Calculate the sector angle for developing a cone of base radius $30\text{ mm}$ and true slant height $90\text{ mm}$.\\[4pt]
\textbf{Answer:}  
$$\theta = \left( \frac{30}{90} \right) \times 360^\circ = \frac{1}{3} \times 360^\circ = \mathbf{120^\circ}$$
\end{vivaqa}

\subsection{Multiple-Choice Questions (MCQs)}

1. Which method of development is used for Pyramids and Cones?  
   A. Parallel Line Method  
   B. Radial Line Method  
   C. Triangulation Method  
   D. Approximate Method  
   \textbf{Answer: B} — Pyramids and Cones are developed using Radial Line Method.

2. All lines on a surface development pattern MUST represent:  
   A. Apparent length  
   B. True length  
   C. Foreshortened length  
   D. Isometric length  
   \textbf{Answer: B} — Every edge on a development pattern must represent True Length ($TL$).

3. Which AutoCAD Boolean command removes one 3D solid from another?  
   A. \texttt{UNION}  
   B. \texttt{SUBTRACT}  
   C. \texttt{INTERSECT}  
   D. \texttt{SLICE}  
   \textbf{Answer: B} — \texttt{SUBTRACT} removes 3D solid volumes.

4. The stretch-out line length for a cylinder of diameter $D$ is:  
   A. $\pi D$  
   B. $2\pi D$  
   C. $\pi D^2$  
   D. $D/2$  
   \textbf{Answer: A} — Stretch-out line length equals perimeter $= \pi D$.

5. Which AutoCAD command automatically generates linked 2D orthographic views from a 3D model?  
   A. \texttt{EXTRUDE}  
   B. \texttt{PRESSPULL}  
   C. \texttt{VIEWBASE}  
   D. \texttt{3DORBIT}  
   \textbf{Answer: C} — \texttt{VIEWBASE} generates 2D views from 3D models.

---

\section{G. Unit 6 Quick Revision Checklist}
\begin{revisionchecklist}
\begin{itemize}
    \item $\text{Parallel Line Method} \to \text{Prisms and Cylinders} \quad (L = \text{Perimeter})$.
    \item $\text{Radial Line Method} \to \text{Pyramids and Cones} \quad \left(\theta = \frac{r}{L} \times 360^\circ\right)$.
    \item $\text{Triangulation Method} \to \text{Transition pieces (square-to-round adapters)}$.
    \item $\text{True Length Rule} \to \text{All edges on development MUST represent True Length } TL$.
    \item $\text{Seam Line} \to \text{Select shortest edge to minimize welding/soldering}$.
    \item $\text{AutoCAD Booleans: } \texttt{UNION (combine)}, \quad \texttt{SUBTRACT (cut hole)}, \quad \texttt{INTERSECT (common volume)}$.
    \item $\texttt{VIEWBASE} \to \text{Auto-generates 2D views from 3D Model Space onto Layout Tab}$.
\end{itemize}
\end{revisionchecklist}
